\documentclass[aspectratio=169]{beamer}

% =========================================================
% TEMPLATE INSTITUCIONAL
% =========================================================
% Compilar con XeLaTeX/LuaLaTeX y -shell-escape si se usan bloques minted.
\input{~/Dropbox/Latex/TemplateSlides/template-slide.tex}

% =========================================================
% CONFIGURACIÓN DEL CURSO
% =========================================================
\renewcommand{\CourseName}{Robótica del Servicio}
\renewcommand{\CourseCode}{ING 01335}
\renewcommand{\Faculty}{Facultad de Ingeniería}
\renewcommand{\Institution}{Politécnico Colombiano Jaime Isaza Cadavid}
\renewcommand{\Instructor}{Deimer Miranda Montoya, MSc.(c).}
\renewcommand{\Email}{deimer\_miranda91162@elpoli.edu.co}
\renewcommand{\InstitutionLogo}{Template/LogoPoli.png}

\title{Robótica del Servicio}
\subtitle{Clase 4: Fundamentos de ROS 2 Humble con Turtlesim}
\author{\Instructor}
\institute{\Institution\\\Faculty}
\date{2026-2}

\begin{document}

\TitleSlide
%\NormasSlide

% =========================================================
\section{Clase 4}

\SectionSlide
{Clase 4: Fundamentos de ROS 2 Humble}
{Del modelo matemático al mensaje robótico}

% =========================================================
\begin{frame}{Propósito de la clase}

    Al finalizar esta sesión, el estudiante estará en capacidad de:

    \begin{itemize}
        \item reconocer la función de ROS 2 dentro de un sistema robótico;
        \item diferenciar nodos, tópicos y mensajes;
        \item ejecutar comandos básicos de ROS 2 desde terminal;
        \item inspeccionar nodos y tópicos activos;
        \item controlar Turtlesim mediante mensajes de velocidad;
        \item observar la pose publicada por el simulador;
        \item relacionar velocidad lineal, velocidad angular y movimiento.
    \end{itemize}

    \vspace{0.3cm}

    \Idea{
        ROS 2 no reemplaza el modelo matemático del robot. ROS 2 permite
        implementar, comunicar y observar los componentes que usan ese modelo.
    }

\end{frame}

% =========================================================
\section{Conexión con las clases anteriores}

\SectionSlide
{Conexión matemática}
{Posición, orientación y movimiento}

% =========================================================
\begin{frame}{¿Qué sabemos hasta ahora?}

    En las clases anteriores trabajamos:

    \begin{itemize}
        \item puntos y vectores en el plano;
        \item distancia euclidiana;
        \item dirección mediante $\operatorname{atan2}$;
        \item coordenadas cartesianas y polares;
        \item matriz de rotación $R(\theta)$;
        \item sistemas de referencia local y global.
    \end{itemize}

    \vspace{0.3cm}

    \Result{
        Ahora comenzaremos a representar computacionalmente estas ideas usando
        ROS 2 y Turtlesim.
    }

\end{frame}

% =========================================================
\begin{frame}{La idea central de la sesión}

    \[
    \text{Modelo matemático}
    \rightarrow
    \text{Mensaje}
    \rightarrow
    \text{Tópico}
    \rightarrow
    \text{Movimiento}
    \]

    \vspace{0.4cm}

    \begin{center}
        \begin{tikzpicture}[
            node distance=1.25cm,
            box/.style={
                draw=BrandBlue,
                thick,
                rounded corners,
                minimum width=2.8cm,
                minimum height=0.9cm,
                align=center
            },
            arrow/.style={
                -{Latex[length=3mm]},
                thick,
                BrandBlue
            }
            ]

            \node[box] (math) {Modelo\\matemático};
            \node[box, right=of math] (msg) {Mensaje\\ROS 2};
            \node[box, right=of msg] (topic) {Tópico};
            \node[box, right=of topic] (move) {Movimiento};

            \draw[arrow] (math) -- (msg);
            \draw[arrow] (msg) -- (topic);
            \draw[arrow] (topic) -- (move);

        \end{tikzpicture}
    \end{center}

    \vspace{0.3cm}

    \Idea{
        El movimiento de un robot puede comenzar como una ecuación, pero en ROS 2
        debe expresarse como un mensaje que viaja por un tópico.
    }

\end{frame}

% =========================================================
\section{¿Qué es ROS 2?}

\SectionSlide
{¿Qué es ROS 2?}
{Infraestructura para sistemas robóticos distribuidos}

% =========================================================
\begin{frame}{¿Qué problema ayuda a resolver ROS 2?}

    Un robot real no suele ser un único programa.

    Normalmente está formado por varios componentes:

    \begin{itemize}
        \item sensores;
        \item actuadores;
        \item algoritmos de control;
        \item estimación de posición;
        \item navegación;
        \item visualización;
        \item registro de datos.
    \end{itemize}

    \vspace{0.3cm}

    \Idea{
        ROS 2 permite que estos componentes se comuniquen de forma organizada.
    }

\end{frame}

% =========================================================
\begin{frame}{Robot como sistema distribuido}

    \begin{center}
        \begin{tikzpicture}[
            node distance=1.0cm,
            box/.style={
                draw=BrandBlue,
                thick,
                rounded corners,
                minimum width=2.8cm,
                minimum height=0.85cm,
                align=center
            },
            smallbox/.style={
                draw=BrandGreen,
                thick,
                rounded corners,
                minimum width=2.4cm,
                minimum height=0.75cm,
                align=center
            },
            arrow/.style={
                -{Latex[length=2.5mm]},
                thick,
                BrandBlue
            }
            ]

            \node[box] (ros) {ROS 2};

            \node[smallbox, above left=of ros] (sensores) {Sensores};
            \node[smallbox, above right=of ros] (control) {Control};
            \node[smallbox, below left=of ros] (nav) {Navegación};
            \node[smallbox, below right=of ros] (viz) {Visualización};

            \draw[arrow] (sensores) -- (ros);
            \draw[arrow] (ros) -- (control);
            \draw[arrow] (nav) -- (ros);
            \draw[arrow] (ros) -- (viz);

        \end{tikzpicture}
    \end{center}

    \vspace{0.3cm}

    \DefBlock[ROS 2]{
        ROS 2 es una infraestructura de software que permite construir sistemas
        robóticos a partir de procesos que se comunican mediante mensajes.
    }

\end{frame}

% =========================================================
\begin{frame}{ROS 2 no es el robot}

    ROS 2 no es:

    \begin{itemize}
        \item un robot físico;
        \item un simulador por sí solo;
        \item un lenguaje de programación;
        \item un controlador automático mágico.
    \end{itemize}

    \vspace{0.3cm}

    ROS 2 es una infraestructura para:

    \begin{itemize}
        \item ejecutar procesos;
        \item comunicar módulos;
        \item transportar datos;
        \item observar el estado del sistema;
        \item organizar software robótico.
    \end{itemize}

\end{frame}

% =========================================================
\begin{frame}{De la matemática al sistema}

    Por ejemplo, en un robot móvil podemos tener:

    \[
    v
    \]

    como velocidad lineal y:

    \[
    \omega
    \]

    como velocidad angular.

    Pero ROS 2 necesita enviar estas cantidades dentro de un mensaje.

    \vspace{0.3cm}

    \[
    (v,\omega)
    \rightarrow
    \texttt{Twist}
    \rightarrow
    \texttt{/cmd\_vel}
    \rightarrow
    \text{movimiento}
    \]

    \vspace{0.3cm}

    \Result{
        La matemática define qué queremos hacer. ROS 2 define cómo lo comunicamos
        entre componentes.
    }

\end{frame}

% =========================================================
\section{Conceptos fundamentales}

\SectionSlide
{Conceptos fundamentales}
{Nodo, tópico y mensaje}

% =========================================================
\begin{frame}{Nodo}

    Un nodo es un proceso que ejecuta una tarea específica dentro de ROS 2.

    En lugar de construir un único programa grande para todo el robot, ROS 2
    permite dividir el sistema en varios nodos especializados.

    \vspace{0.3cm}

    \DefBlock[Nodo]{
        Un nodo es una unidad de ejecución. Cada nodo se encarga de una parte
        concreta del sistema robótico y puede comunicarse con otros nodos.
    }

\end{frame}

% =========================================================
\begin{frame}{Nodo: idea robótica}

    En un robot móvil podríamos tener:

    \begin{itemize}
        \item un nodo que lee sensores de distancia;
        \item un nodo que estima la posición del robot;
        \item un nodo que calcula una acción de control;
        \item un nodo que envía velocidades a los motores;
        \item un nodo que registra datos para análisis posterior.
    \end{itemize}

    \vspace{0.3cm}

    \Idea{
        Cada nodo tiene una responsabilidad clara. Esto facilita probar,
        corregir y reemplazar partes del sistema sin modificar todo el robot.
    }

\end{frame}

% =========================================================
\begin{frame}{Nodo: separación de responsabilidades}

    \begin{center}
        \begin{tikzpicture}[
            node distance=1.0cm,
            nodebox/.style={
                ellipse,
                draw=BrandBlue,
                thick,
                minimum width=3.0cm,
                minimum height=1.0cm,
                align=center
            },
            arrow/.style={
                -{Latex[length=3mm]},
                thick,
                BrandBlue
            }
            ]

            \node[nodebox] (sens) {Lectura\\sensores};
            \node[nodebox, right=of sens] (est) {Estimación\\pose};
            \node[nodebox, right=of est] (ctrl) {Control\\movimiento};

            \draw[arrow] (sens) -- (est);
            \draw[arrow] (est) -- (ctrl);

        \end{tikzpicture}
    \end{center}

    \vspace{0.3cm}

    \Result{
        Un sistema robótico se vuelve más organizado cuando cada nodo cumple
        una función específica.
    }

\end{frame}

% =========================================================
\begin{frame}{Tópico}

    Un tópico es un canal de comunicación por donde viajan mensajes.

    Los nodos pueden:

    \begin{itemize}
        \item publicar información en un tópico;
        \item suscribirse a un tópico para recibir información.
    \end{itemize}

    \vspace{0.3cm}

    \DefBlock[Tópico]{
        Un tópico no ejecuta cálculos. Su función es transportar mensajes
        entre nodos que producen información y nodos que la necesitan.
    }

\end{frame}

% =========================================================
\begin{frame}{Tópico: analogía}

    Puede pensarse en un tópico como un canal con un nombre.

    Por ejemplo:

    \[
    \texttt{/turtle1/cmd\_vel}
    \]

    es el canal por donde se envían comandos de velocidad.

    \[
    \texttt{/turtle1/pose}
    \]

    es el canal por donde Turtlesim publica posición y orientación.

    \vspace{0.3cm}

    \Idea{
        En ROS 2 no basta con tener datos. Es necesario saber en qué tópico
        se publican y qué tipo de mensaje transportan.
    }

\end{frame}

% =========================================================
\begin{frame}{Comunicación por tópicos}

    \begin{center}
        \begin{tikzpicture}[
            node distance=1.8cm,
            nodebox/.style={
                ellipse,
                draw=BrandBlue,
                thick,
                minimum width=3.0cm,
                minimum height=1.05cm,
                align=center
            },
            topicbox/.style={
                rectangle,
                draw=BrandYellow,
                thick,
                minimum width=3.0cm,
                minimum height=0.9cm,
                align=center
            },
            arrow/.style={
                -{Latex[length=3mm]},
                thick,
                BrandBlue
            }
            ]

            \node[nodebox] (pub) {Nodo\\publicador};
            \node[topicbox, right=of pub] (topic) {Tópico};
            \node[nodebox, right=of topic] (sub) {Nodo\\suscriptor};

            \draw[arrow] (pub) -- node[above] {mensaje} (topic);
            \draw[arrow] (topic) -- node[above] {mensaje} (sub);

        \end{tikzpicture}
    \end{center}

    \vspace{0.3cm}

    \Idea{
        El publicador no necesita saber quién recibe el mensaje. El suscriptor
        solo necesita conocer el tópico al cual escuchar.
    }

\end{frame}

% =========================================================
\begin{frame}{Varios nodos pueden usar el mismo tópico}

    Un mismo tópico puede tener:

    \begin{itemize}
        \item varios publicadores;
        \item varios suscriptores;
        \item un publicador y múltiples consumidores de información.
    \end{itemize}

    \vspace{0.2cm}

    \begin{center}
        \begin{tikzpicture}[
            node distance=0.9cm,
            nodebox/.style={
                ellipse,
                draw=BrandBlue,
                thick,
                minimum width=2.7cm,
                minimum height=0.9cm,
                align=center
            },
            topicbox/.style={
                rectangle,
                draw=BrandYellow,
                thick,
                minimum width=3.2cm,
                minimum height=0.9cm,
                align=center
            },
            arrow/.style={
                -{Latex[length=2.6mm]},
                thick,
                BrandBlue
            }
            ]

            \node[topicbox] (topic) {\texttt{/turtle1/pose}};

            \node[nodebox, above left=of topic] (sim) {Turtlesim};
            \node[nodebox, below left=of topic] (log) {Registro\\datos};
            \node[nodebox, above right=of topic] (echo) {Terminal\\echo};
            \node[nodebox, below right=of topic] (plot) {Gráfica\\Python};

            \draw[arrow] (sim) -- (topic);
            \draw[arrow] (topic) -- (echo);
            \draw[arrow] (topic) -- (plot);
            \draw[arrow] (topic) -- (log);

        \end{tikzpicture}
    \end{center}

\end{frame}

% =========================================================
\begin{frame}{Mensaje}

    Un mensaje define la estructura de los datos que viajan por un tópico.

    No todos los tópicos transportan el mismo tipo de información.

    \vspace{0.3cm}

    Por ejemplo:

    \[
    \texttt{/turtle1/cmd\_vel}
    \]

    transporta un mensaje de velocidad.

    \[
    \texttt{/turtle1/pose}
    \]

    transporta un mensaje de posición y orientación.

    \vspace{0.3cm}

    \DefBlock[Mensaje]{
        Un mensaje es el formato de datos que ROS 2 utiliza para comunicar
        nodos. Define qué campos existen y qué tipo de dato tiene cada campo.
    }

\end{frame}

% =========================================================
\begin{frame}{Mensaje: estructura}

    Un mensaje no es solo un número. Puede contener varios campos.

    Por ejemplo, un mensaje de velocidad tipo \texttt{Twist} contiene:

    \[
    \texttt{linear}
    \qquad
    \text{y}
    \qquad
    \texttt{angular}
    \]

    \vspace{0.2cm}

    Cada parte puede tener componentes en:

    \[
    x,\quad y,\quad z
    \]

    \vspace{0.3cm}

    \Idea{
        Para mover Turtlesim en el plano usaremos principalmente
        \texttt{linear.x} y \texttt{angular.z}.
    }

\end{frame}

% =========================================================
\begin{frame}{Mensaje: relación con la matemática}

    En robótica móvil plana usamos:

    \[
    v
    \]

    para velocidad lineal, y:

    \[
    \omega
    \]

    para velocidad angular.

    En ROS 2 estas variables se escriben dentro del mensaje:

    \[
    v \rightarrow \texttt{linear.x}
    \]

    \[
    \omega \rightarrow \texttt{angular.z}
    \]

    \vspace{-0.3cm}

    \Result{
        La matemática define la variable. El mensaje define cómo esa variable
        viaja dentro del sistema ROS 2.
    }

\end{frame}

% =========================================================
\begin{frame}{Publicador}

    Un publicador es la parte de un nodo que envía mensajes a un tópico.

    Ejemplo:

    \[
    \texttt{/teleop\_turtle}
    \]

    publica comandos de velocidad en:

    \[
    \texttt{/turtle1/cmd\_vel}
    \]

    \vspace{0.3cm}

    \DefBlock[Publicador]{
        Un publicador produce mensajes y los coloca en un tópico específico.
    }

\end{frame}

% =========================================================
\begin{frame}{Suscriptor}

    Un suscriptor es la parte de un nodo que recibe mensajes desde un tópico.

    Ejemplo:

    \[
    \texttt{/turtlesim}
    \]

    se suscribe al tópico:

    \[
    \texttt{/turtle1/cmd\_vel}
    \]

    para recibir velocidades y mover la tortuga.

    \vspace{0.3cm}

    \DefBlock[Suscriptor]{
        Un suscriptor escucha un tópico y ejecuta una acción cuando recibe
        un mensaje.
    }

\end{frame}

% =========================================================
\begin{frame}{Publicador y suscriptor}

    \begin{columns}[T]

        \begin{column}{0.48\textwidth}

            \textbf{Publicador}

            \begin{itemize}
                \item genera datos;
                \item los envía a un tópico;
                \item no controla directamente quién los recibe.
            \end{itemize}

            \[
            \texttt{publish()}
            \]

        \end{column}

        \begin{column}{0.48\textwidth}

            \textbf{Suscriptor}

            \begin{itemize}
                \item escucha un tópico;
                \item recibe mensajes;
                \item procesa la información recibida.
            \end{itemize}

            \[
            \texttt{subscribe()}
            \]

        \end{column}

    \end{columns}

    \vspace{0.3cm}

    \Result{
        Este patrón será la base de casi todo el trabajo con ROS 2.
    }

\end{frame}

% =========================================================
\begin{frame}{Ejemplo completo en Turtlesim}

    \begin{center}
        \begin{tikzpicture}[
            node distance=1.2cm,
            nodebox/.style={
                ellipse,
                draw=BrandBlue,
                thick,
                minimum width=3.0cm,
                minimum height=1.0cm,
                align=center
            },
            topicbox/.style={
                rectangle,
                draw=BrandYellow,
                thick,
                minimum width=3.4cm,
                minimum height=0.9cm,
                align=center
            },
            arrow/.style={
                -{Latex[length=2.8mm]},
                thick,
                BrandBlue
            }
            ]

            \node[nodebox] (teleop) {\texttt{/teleop\_turtle}};
            \node[topicbox, right=of teleop] (cmd) {\texttt{/turtle1/cmd\_vel}};
            \node[nodebox, right=of cmd] (sim) {\texttt{/turtlesim}};

            \node[topicbox, below=of sim] (pose) {\texttt{/turtle1/pose}};
            \node[nodebox, left=of pose] (echo) {\texttt{ros2 topic echo}};

            \draw[arrow] (teleop) -- node[above] {\texttt{Twist}} (cmd);
            \draw[arrow] (cmd) -- node[above] {\texttt{Twist}} (sim);
            \draw[arrow] (sim) -- node[right] {\texttt{Pose}} (pose);
            \draw[arrow] (pose) -- node[below] {\texttt{Pose}} (echo);

        \end{tikzpicture}
    \end{center}

\end{frame}

% =========================================================
\begin{frame}{Lectura del diagrama}

    En el diagrama anterior:

    \begin{itemize}
        \item \texttt{/teleop\_turtle} es un nodo;
        \item \texttt{/turtle1/cmd\_vel} es un tópico;
        \item el mensaje que viaja allí es \texttt{Twist};
        \item \texttt{/turtlesim} recibe ese mensaje;
        \item \texttt{/turtlesim} publica la pose en \texttt{/turtle1/pose};
        \item la terminal puede suscribirse para observar esa pose.
    \end{itemize}

    \vspace{0.3cm}

    \Idea{
        En ROS 2 siempre debemos preguntarnos:
        ¿quién publica?, ¿en qué tópico?, ¿qué tipo de mensaje viaja?,
        ¿quién se suscribe?
    }

\end{frame}

% =========================================================
\section{Exploración de mensajes desde terminal}

\SectionSlide
{Exploración de mensajes ROS 2}
{De un dato simple a una estructura de velocidad}

% =========================================================
\begin{frame}{Antes de usar Turtlesim}

    Ya sabemos que un mensaje define la estructura de los datos que viajan
    entre nodos.

    Pero antes de utilizar mensajes en un robot, podemos explorarlos
    directamente desde la terminal.

    \vspace{0.3cm}

    Comenzaremos con estructuras simples:

    \[
    \texttt{Int32}
    \rightarrow
    \texttt{Float64}
    \]

    y progresivamente llegaremos a:

    \[
    \texttt{Vector3}
   \rightarrow
    \boxed{\texttt{Twist}}
    \]

    \Idea{
        La complejidad de un mensaje depende de la información que necesitamos
        representar.
    }

\end{frame}

% =========================================================
\begin{frame}{Nombre completo de un tipo de mensaje}

    En ROS 2 encontraremos nombres como:

    \[
    \boxed{
        \texttt{std\_msgs/msg/Int32}
    }
    \]

    Este nombre contiene tres partes:

    \[
    \underbrace{\texttt{std\_msgs}}_{\text{paquete}}
    /
    \underbrace{\texttt{msg}}_{\text{tipo de interfaz}}
    /
    \underbrace{\texttt{Int32}}_{\text{mensaje}}
    \]

    \vspace{0.3cm}

    \DefBlock[Notación]{
        La expresión \texttt{paquete/msg/Mensaje} identifica de manera
        completa un tipo de mensaje dentro de ROS 2.
    }

\end{frame}

% =========================================================
\begin{frame}{¿Qué es \texorpdfstring{\texttt{std\_msgs}}{std msgs}?}

    El paquete:

    \[
    \texttt{std\_msgs}
    \]

    contiene mensajes básicos utilizados para transportar valores sencillos.

    Por ejemplo:

    \begin{itemize}
        \item números enteros;
        \item números reales;
        \item valores booleanos;
        \item cadenas de caracteres.
    \end{itemize}

    \vspace{0.3cm}

    \Idea{
        Aunque estos mensajes son simples, permiten comprender cómo ROS 2
        encapsula un dato antes de enviarlo.
    }

\end{frame}

% =========================================================
\begin{frame}[fragile]{Explorar las interfaces disponibles}

    Podemos consultar las interfaces conocidas por ROS 2:

    \begin{bashcodedark}
        ros2 interface list
    \end{bashcodedark}

    La lista puede ser extensa.

    Para buscar mensajes relacionados con \texttt{std\_msgs}:

    \begin{bashcodedark}
        ros2 interface list | grep std_msgs
    \end{bashcodedark}

    \vspace{0.3cm}

    \DefBlock[ros2 interface list]{
        Muestra las interfaces disponibles en la instalación actual de ROS 2.
    }

\end{frame}

% =========================================================
\begin{frame}[fragile]{Explorar un paquete de interfaces}

    También podemos consultar las interfaces pertenecientes a un paquete:

    \begin{bashcodedark}
        ros2 interface package std_msgs
    \end{bashcodedark}

    Entre los resultados podremos encontrar mensajes como:

    \begin{consolecodedark}
        std_msgs/msg/Bool
        std_msgs/msg/Float64
        std_msgs/msg/Int32
        std_msgs/msg/String
    \end{consolecodedark}

    \vspace{0.3cm}

    \Idea{
        Un paquete puede contener varios tipos de mensajes diseñados para
        representar diferentes clases de información.
    }

\end{frame}

% =========================================================
\section{Mensajes simples}

\SectionSlide
{Mensajes simples}
{Un dato también necesita una estructura}

% =========================================================
\begin{frame}[fragile]{Mensaje Int32}

    Consultemos:

    \begin{bashcodedark}
        ros2 interface show std_msgs/msg/Int32
    \end{bashcodedark}

    La estructura es:

    \begin{consolecodedark}
        int32 data
    \end{consolecodedark}

    Esto significa que el mensaje contiene un campo llamado:

    \[
    \texttt{data}
    \]

    de tipo:

    \[
    \texttt{int32}
    \]

\end{frame}

% =========================================================
\begin{frame}{Interpretación de Int32}

    Podemos imaginar el mensaje como un pequeño contenedor:

    \begin{center}
        \begin{tikzpicture}[
            box/.style={
                rectangle,
                draw=BrandYellow,
                thick,
                minimum width=5.0cm,
                minimum height=1.1cm,
                align=center
            }
            ]

            \node[box] {
                \texttt{std\_msgs/msg/Int32}\\[1mm]
                \texttt{data : int32}
            };

        \end{tikzpicture}
    \end{center}

    Por ejemplo:

    \[
    \texttt{data}=25
    \]

    \vspace{0.3cm}

    \Idea{
        El valor 25 no viaja ``solo''. Viaja dentro de una estructura
        definida por el mensaje \texttt{Int32}.
    }

\end{frame}

% =========================================================
\begin{frame}{Tipo del campo y tipo del mensaje}

    Es importante diferenciar:

    \[
    \boxed{\texttt{int32}}
    \]

    de:

    \[
    \boxed{\texttt{std\_msgs/msg/Int32}}
    \]

    El primero es un:

    \[
    \text{tipo de dato}
    \]

    El segundo es un:

    \[
    \text{tipo de mensaje ROS 2}
    \]

    \vspace{-0.3cm}

    \Warning{
        \texttt{int32} no es un tópico ni un mensaje completo.
        Es el tipo de dato utilizado por el campo \texttt{data}.
    }

\end{frame}

% =========================================================
\begin{frame}[fragile]{Mensaje Float64}

    Ahora consulte:

    \begin{bashcodedark}
        ros2 interface show std_msgs/msg/Float64
    \end{bashcodedark}

    La estructura es:

    \begin{consolecodedark}
        float64 data
    \end{consolecodedark}

    \vspace{0.3cm}

    Un posible contenido sería:

    \[
    \texttt{data}=3.1416
    \]

    \DefBlock[Float64]{
        \texttt{std\_msgs/msg/Float64} permite transportar un número real
        de precisión doble dentro del campo \texttt{data}.
    }

\end{frame}

% =========================================================
\begin{frame}{Ejemplo robótico con Float64}

    Supongamos que un sensor mide una distancia:

    \[
    d=0.73~\text{m}
    \]

    El valor matemático es:

    \[
    0.73
    \]

    Una posible representación mediante un mensaje simple sería:

    \[
    \boxed{
        \texttt{data: 0.73}
    }
    \]

    \vspace{0.3cm}

    \begin{center}
        \begin{tikzpicture}[
            box/.style={
                rectangle,
                draw=BrandYellow,
                thick,
                minimum width=4.5cm,
                minimum height=1.2cm,
                align=center
            }
            ]

            \node[box] {
                \texttt{Float64}\\
                \texttt{data = 0.73}
            };

        \end{tikzpicture}
    \end{center}

\end{frame}

% =========================================================
\begin{frame}{Pero un robot necesita estructuras más grandes}

    Un único número puede representar:

    \[
    d
     \rightarrow
    v
    \rightarrow
    \theta
    \]

    Pero algunas cantidades necesitan varias componentes.

    Por ejemplo, un vector tridimensional:

    \[
    \mathbf{v}
    =
    \begin{bmatrix}
        v_x,
        v_y,
        v_z
    \end{bmatrix}^T
    \]

    requiere almacenar tres valores.

    \vspace{0.3cm}

    \Idea{
        Cuando la información tiene varias componentes, necesitamos mensajes
        con estructuras más complejas.
    }

\end{frame}

% =========================================================
\section{Mensajes compuestos}

\SectionSlide
{Mensajes compuestos}
{Cuando un solo campo ya no es suficiente}

% =========================================================
\begin{frame}[fragile]{Explorar Vector3}

    Consultemos:

    \begin{bashcodedark}
        ros2 interface show geometry_msgs/msg/Vector3
    \end{bashcodedark}

    Su estructura es:

    \begin{consolecodedark}
        float64 x
        float64 y
        float64 z
    \end{consolecodedark}

    \vspace{0.3cm}

    Ahora el mensaje contiene:

    \[
    3
    \]

    campos independientes.

\end{frame}

% =========================================================
\begin{frame}{Interpretación de Vector3}

    Podemos representarlo como:

    \[
    \mathbf{v}
    =
    \begin{bmatrix}
        x\\
        y\\
        z
    \end{bmatrix}
    \]

    En ROS 2:

    \begin{center}
        \begin{tikzpicture}[
            box/.style={
                rectangle,
                draw=BrandYellow,
                thick,
                minimum width=5.0cm,
                minimum height=0.75cm,
                align=left
            }
            ]

            \node[box] (title) {\texttt{geometry\_msgs/msg/Vector3}};
            \node[box, below=0cm of title] (x) {\texttt{x : float64}};
            \node[box, below=0cm of x] (y) {\texttt{y : float64}};
            \node[box, below=0cm of y] (z) {\texttt{z : float64}};

        \end{tikzpicture}
    \end{center}

\end{frame}

% =========================================================
\begin{frame}{Ejemplo de Vector3}

    Supongamos:

    \[
    \mathbf{v}
    =
    \begin{bmatrix}
        2.0\\
        -1.5\\
        0.4
    \end{bmatrix}
    \]

    El mensaje contendría:

    \[
    \texttt{x: 2.0}
    \]

    \[
    \texttt{y: -1.5}
    \]

    \[
    \texttt{z: 0.4}
    \]

    \vspace{-0.3cm}

    \Result{
        El mensaje \texttt{Vector3} permite agrupar tres componentes relacionadas
        dentro de una sola estructura.
    }

\end{frame}

% =========================================================
\begin{frame}{La estructura puede contener otras estructuras}

    Hasta ahora tenemos:

    \[
    \texttt{Int32}
    \rightarrow
    \text{un campo}
    \]

    \[
    \texttt{Float64}
    \rightarrow
    \text{un campo}
    \]

    \[
    \texttt{Vector3}
    \rightarrow
    \text{tres campos}
    \]

    Pero ROS 2 también permite que un mensaje contenga:

    \[
    \boxed{\text{otros mensajes}}
    \]

    \vspace{0.3cm}

    \Idea{
        Esta composición permite construir estructuras de datos progresivamente
        más completas.
    }

\end{frame}

% =========================================================
\section{Llegamos al mensaje Twist}

\SectionSlide
{Mensaje Twist}
{Dos vectores dentro de un mensaje}

% =========================================================
\begin{frame}[fragile]{Explorar Twist desde terminal}

    Ejecute:

    \begin{bashcodedark}
        ros2 interface show geometry_msgs/msg/Twist
    \end{bashcodedark}

    Su estructura conceptual es:

    \begin{consolecodedark}
        Vector3 linear
        Vector3 angular
    \end{consolecodedark}

    \vspace{0.3cm}

    Esto significa que \texttt{Twist} contiene:

    \[
    2
    \]

    mensajes de tipo:

    \[
    \texttt{Vector3}
    \]

\end{frame}

% =========================================================
\begin{frame}{Construcción jerárquica de Twist}

    \begin{center}
        \begin{tikzpicture}[
            node distance=0.8cm,
            box/.style={
                rectangle,
                draw=BrandYellow,
                thick,
                minimum width=3.2cm,
                minimum height=0.85cm,
                align=center
            },
            arrow/.style={
                -{Latex[length=2.6mm]},
                thick,
                BrandBlue
            }
            ]

            \node[box] (twist) {\texttt{Twist}};

            \node[box, below left=of twist] (linear) {\texttt{linear}\\\texttt{Vector3}};
            \node[box, below right=of twist] (angular) {\texttt{angular}\\\texttt{Vector3}};

            \node[box, below=of linear] (lxyz) {
                \texttt{x}\\
                \texttt{y}\\
                \texttt{z}
            };

            \node[box, below=of angular] (axyz) {
                \texttt{x}\\
                \texttt{y}\\
                \texttt{z}
            };

            \draw[arrow] (twist) -- (linear);
            \draw[arrow] (twist) -- (angular);

            \draw[arrow] (linear) -- (lxyz);
            \draw[arrow] (angular) -- (axyz);

        \end{tikzpicture}
    \end{center}

\end{frame}

% =========================================================
\begin{frame}{Estructura completa de Twist}

    Podemos pensar en el mensaje como:

    \[
    \texttt{Twist}
    =
    \left\{
    \begin{array}{l}
        \texttt{linear}=
        \left\{
        \begin{array}{l}
            x\\
            y\\
            z
        \end{array}
        \right.
        \\[4mm]
        \texttt{angular}=
        \left\{
        \begin{array}{l}
            x\\
            y\\
            z
        \end{array}
        \right.
    \end{array}
    \right.
    \]

    \vspace{0.3cm}

    Por tanto, existen seis componentes numéricas:

    \[
    v_x,\;v_y,\;v_z,\;
    \omega_x,\;\omega_y,\;\omega_z
    \]

\end{frame}

% =========================================================
\begin{frame}{Interpretación física}

    El bloque:

    \[
    \texttt{linear}
    \]

    representa velocidad lineal:

    \[
    \mathbf{v}
    =
    \begin{bmatrix}
        v_x\\
        v_y\\
        v_z
    \end{bmatrix}
    \]

    Mientras que:

    \[
    \texttt{angular}
    \]

    representa velocidad angular:

    \[
    \boldsymbol{\omega}
    =
    \begin{bmatrix}
        \omega_x\\
        \omega_y\\
        \omega_z
    \end{bmatrix}
    \]

\end{frame}

% =========================================================
\begin{frame}{¿Por qué usamos \texorpdfstring{\texttt{linear.x}}{linear.x}?}

    Para un robot móvil que se desplaza sobre un plano, consideraremos:

    \[
    X
    \]

    como la dirección hacia adelante.

    Por tanto:

    \[
    \boxed{
        v=\texttt{linear.x}
    }
    \]

    \vspace{0.3cm}

    \begin{center}
        \begin{tikzpicture}[scale=1.0]

            \coordinate (R) at (2.5,1.5);

            \begin{scope}[rotate around={-90:(0.0,0.0)}]
            % Cuerpo circular del robot
            \draw[
            fill=BrandBlue!15,
            draw=BrandBlue,
            thick
            ]
            (R) circle (0.75);

            % Rueda izquierda
            \draw[
            fill=black!70,
            draw=black,
            thick,
            rounded corners=1pt
            ]
            ($(R)+(-0.9,-0.32)$) rectangle ($(R)+(-0.68,0.32)$);

            % Rueda derecha
            \draw[
            fill=black!70,
            draw=black,
            thick,
            rounded corners=1pt
            ]
            ($(R)+(0.68,-0.32)$) rectangle ($(R)+(0.9,0.32)$);

            % Indicador del frente del robot
            \draw[
            -{Latex[length=2.5mm]},
            very thick,
            BrandBlue
            ]
            (R) -- ++(0.0,0.55);

            % Vector de velocidad longitudinal
            \draw[
            -{Latex[length=3mm]},
            very thick,
            BrandGreen
            ]
            ($(R)+(0,0.8)$) -- ++(0,1.4)
            node[above] {$v_x$};
            \end{scope}

            % Etiqueta
            \node[below=1cm] at (R) {robot};

        \end{tikzpicture}
    \end{center}

\end{frame}

% =========================================================
\begin{frame}{¿Por qué usamos \texorpdfstring{\texttt{angular.z}}{angular.z}?}

    El robot se mueve sobre el plano:

    \[
    XY
    \]

    La rotación que cambia su orientación ocurre alrededor del eje:

    \[
    Z
    \]

    que es perpendicular al plano.

    Por tanto:

    \[
    \boxed{
        \omega=\texttt{angular.z}
    }
    \]

\end{frame}

% =========================================================
\begin{frame}{Velocidad angular alrededor de Z}

\begin{center}
    \begin{tikzpicture}[scale=0.95]

        % Ejes globales
        \draw[-{Latex[length=3mm]}, thick]
        (0,0) -- (5,0)
        node[right] {$X$};

        \draw[-{Latex[length=3mm]}, thick]
        (0,0) -- (0,3.8)
        node[above] {$Y$};

        % Posición del robot
        \coordinate (R) at (2.3,1.6);

        % Robot visto desde arriba
        \draw[
        fill=BrandBlue!15,
        draw=BrandBlue,
        thick
        ]
        (R) circle (0.55);

        % Rueda izquierda
        \draw[
        fill=black!70,
        draw=black,
        thick,
        rounded corners=1pt
        ]
        ($(R)+(-0.68,-0.23)$)
        rectangle
        ($(R)+(-0.53,0.23)$);

        % Rueda derecha
        \draw[
        fill=black!70,
        draw=black,
        thick,
        rounded corners=1pt
        ]
        ($(R)+(0.53,-0.23)$)
        rectangle
        ($(R)+(0.68,0.23)$);

        % Frente del robot
        \draw[
        -{Latex[length=2.5mm]},
        very thick,
        BrandBlue
        ]
        (R) -- ++(0,0.42);

        % Velocidad lineal
        \draw[
        -{Latex[length=3mm]},
        very thick,
        BrandGreen
        ]
        ($(R)+(0,0.58)$) -- ++(0,1.1)
        node[above] {$v_x$};

        % Velocidad angular alrededor de Z
        \draw[
        -{Latex[length=2.5mm]},
        BrandYellow,
        very thick
        ]
        ($(R)+(0.8,0)$)
        arc[
        start angle=0,
        end angle=120,
        radius=0.8
        ];

        \node[BrandYellow]
        at ($(R)+(0.95,0.75)$)
        {$\omega_z$};

        % Eje Z saliendo del plano
        \draw[BrandBlue, thick]
        (R) circle (0.11);

        \fill[BrandBlue]
        (R) circle (1.7pt);

%        \node[
%        BrandBlue,
%        below right=2pt
%        ] at (R) {$Z$};

        % Etiqueta opcional
        \node[below=0.8cm] at (R) {robot};

    \end{tikzpicture}
\end{center}

    \Idea{
        Aunque el movimiento ocurre en el plano $XY$, la rotación se describe
        alrededor del eje perpendicular $Z$.
    }

\end{frame}

% =========================================================
\begin{frame}{Del mensaje al movimiento planar}

    Para nuestro caso:

    \[
        \begin{matrix}
        \texttt{linear.x}=v \\
        \texttt{linear.y}=0\\
        \texttt{linear.z}=0
    \end{matrix}
    \]


    y:

    \[
        \begin{matrix}
    \texttt{angular.x}=0\\
    \texttt{angular.y}=0\\
    \texttt{angular.z}=\omega
    \end{matrix}
    \]

    Por tanto:

    \[
    \boxed{
        \mathbf{u}
        =
        \begin{bmatrix}
            v\\
            \omega
        \end{bmatrix}
    }
    \]

\end{frame}

% =========================================================
\begin{frame}{De la matemática a Twist}

    \begin{center}
        \begin{tikzpicture}[
            node distance=1.0cm,
            box/.style={
                rectangle,
                draw=BrandYellow,
                thick,
                minimum width=3.0cm,
                minimum height=0.85cm,
                align=center
            },
            arrow/.style={
                -{Latex[length=3mm]},
                thick,
                BrandBlue
            }
            ]

            \node[box] (math) {
                $\displaystyle
                \mathbf{u}=
                \begin{bmatrix}
                    v\\
                    \omega
                \end{bmatrix}$
            };

            \node[box, right=of math] (twist) {\texttt{Twist}};

            \node[box, right=of twist] (fields) {
                \texttt{linear.x = v}\\
                \texttt{angular.z = $\omega$}
            };

            \draw[arrow] (math) -- (twist);
            \draw[arrow] (twist) -- (fields);

        \end{tikzpicture}
    \end{center}

    \vspace{0.3cm}

    \Result{
        Ahora el mensaje \texttt{Twist} deja de ser una estructura abstracta:
        es una representación computacional de variables físicas de movimiento.
    }

\end{frame}

% =========================================================
\section{Ejercicio de exploración}

\SectionSlide
{Exploración guiada}
{Leer una interfaz antes de utilizarla}

% =========================================================
\begin{frame}[fragile]{Secuencia de exploración}

    Ejecute en orden:

    \begin{bashcodedark}
        ros2 interface show std_msgs/msg/Int32
    \end{bashcodedark}

    \begin{bashcodedark}
        ros2 interface show std_msgs/msg/Float64
    \end{bashcodedark}

    \begin{bashcodedark}
        ros2 interface show geometry_msgs/msg/Vector3
    \end{bashcodedark}

    \begin{bashcodedark}
        ros2 interface show geometry_msgs/msg/Twist
    \end{bashcodedark}

\end{frame}

% =========================================================
\begin{frame}{Compare las estructuras}

    Complete mentalmente la progresión:

    \[
    \begin{matrix}
    \boxed{
        \texttt{Int32}
    }
    \\
    \downarrow
    \\
    \boxed{
        \texttt{Float64}
    }
    \\
    \downarrow
    \\
    \boxed{
        \texttt{Vector3}
    }
    \\
    \downarrow
    \\
    \boxed{
        \texttt{Twist}
    }
\end{matrix}
    \]

    \vspace{0.3cm}

    \Idea{
        No estamos cambiando simplemente de tipo numérico. Estamos pasando
        de estructuras simples a estructuras compuestas capaces de representar
        información física más rica.
    }

\end{frame}

% =========================================================
\begin{frame}{Preguntas antes de continuar}

    Analice:

    \begin{enumerate}
        \item ¿Cuál es la diferencia entre \texttt{float64} y
        \texttt{std\_msgs/msg/Float64}?

        \item ¿Cuántos campos posee \texttt{Vector3}?

        \item ¿Qué tipo de mensaje son \texttt{linear} y \texttt{angular}
        dentro de \texttt{Twist}?

        \item ¿Cuántos valores numéricos contiene finalmente un \texttt{Twist}?

        \item ¿Por qué para movimiento planar utilizaremos principalmente
        \texttt{linear.x} y \texttt{angular.z}?
    \end{enumerate}

\end{frame}

% =========================================================
\begin{frame}{Estamos listos para utilizar Turtlesim}

    Ahora sabemos interpretar:

    \[
    \texttt{geometry\_msgs/msg/Twist}
    \]

    y entendemos que contiene:

    \[
    \boxed{
        v \rightarrow \texttt{linear.x}
    }
    \]

    \[
    \boxed{
        \omega \rightarrow \texttt{angular.z}
    }
    \]

    \vspace{0.3cm}

    \Result{
        El siguiente paso será identificar qué nodo publica este mensaje,
        por qué tópico viaja y qué nodo lo recibe para producir movimiento.
    }

\end{frame}

% =========================================================
\section{Turtlesim}

\SectionSlide
{Turtlesim}
{Un laboratorio mínimo para entender ROS 2}

% =========================================================
\begin{frame}{¿Por qué usar Turtlesim?}

Turtlesim permite estudiar ROS 2 sin depender todavía del robot físico.

Con Turtlesim podemos practicar:

\begin{itemize}
    \item ejecución de nodos;
    \item publicación de mensajes;
    \item suscripción a tópicos;
    \item comandos de inspección;
    \item movimiento en el plano;
    \item interpretación de pose.
\end{itemize}

\vspace{0.3cm}

\Idea{
    Turtlesim es simple, pero contiene las ideas esenciales que luego veremos
    en un robot móvil real.
}

\end{frame}

% =========================================================
\begin{frame}{Arquitectura básica con Turtlesim}

\begin{center}
    \begin{tikzpicture}[
        node distance=1.1cm,
        box/.style={
            ellipse,
            draw=BrandBlue,
            thick,
            rounded corners,
            minimum width=3.2cm,
            minimum height=0.9cm,
            align=center
        },
        topic/.style={
            draw=BrandYellow,
            thick,
            rounded corners,
            minimum width=3.0cm,
            minimum height=0.8cm,
            align=center
        },
        arrow/.style={
            -{Latex[length=2.8mm]},
            thick,
            BrandBlue
        }
        ]

        \node[box] (teleop) {Nodo\\teleop};
        \node[topic, right=of teleop] (cmd) {\texttt{/turtle1/cmd\_vel}};
        \node[box, right=of cmd] (turtle) {Nodo\\turtlesim};

        \node[topic, below=of turtle] (pose) {\texttt{/turtle1/pose}};
        \node[box, left=of pose] (echo) {Terminal\\\texttt{echo}};

        \draw[arrow] (teleop) -- (cmd);
        \draw[arrow] (cmd) -- (turtle);
        \draw[arrow] (turtle) -- (pose);
        \draw[arrow] (pose) -- (echo);

    \end{tikzpicture}
\end{center}

\end{frame}

% =========================================================
\begin{frame}{Movimiento en Turtlesim}

Turtlesim se mueve a partir de comandos de velocidad.

La tortuga recibe mensajes en:

\[
\texttt{/turtle1/cmd\_vel}
\]

y publica su pose en:

\[
\texttt{/turtle1/pose}
\]

\vspace{0.3cm}

\[
\text{velocidad}
\rightarrow
\text{movimiento}
\rightarrow
\text{pose}
\]

\vspace{0.3cm}

\Result{
    Esta relación es la primera conexión práctica entre modelo de movimiento
    y comunicación ROS 2.
}

\end{frame}

% =========================================================
\section{Preparar la terminal}

\SectionSlide
{Preparar la terminal}
{Cargar ROS 2 Humble}

% =========================================================
\begin{frame}[fragile]{Verificar ROS 2}

Abra una terminal y verifique la distribución activa:

\begin{bashcodedark}
    echo $ROS_DISTRO
\end{bashcodedark}

La salida esperada es:

\begin{bashcode}
    humble
\end{bashcode}

\vspace{0.3cm}

\Warning{
    Si no aparece \texttt{humble}, debe cargar el entorno de ROS 2.
}

\end{frame}

% =========================================================
\begin{frame}[fragile]{Cargar ROS 2 Humble}

Ejecute:

\begin{bashcodedark}
    source /opt/ros/humble/setup.bash
\end{bashcodedark}

Verifique nuevamente:

\begin{bashcodedark}
    echo $ROS_DISTRO
\end{bashcodedark}

La salida debe ser:

\begin{bashcode}
    humble
\end{bashcode}

\vspace{0.3cm}

\Idea{
    El comando \texttt{source} carga las variables necesarias para que la
    terminal reconozca los comandos y paquetes de ROS 2.
}

\end{frame}

% =========================================================
\section{Ejecutar Turtlesim}

\SectionSlide
{Ejecutar Turtlesim}
{Primeros nodos ROS 2}

% =========================================================
\begin{frame}[fragile]{Instalar Turtlesim}

Si Turtlesim no se encuentra instalado, ejecute:

\begin{bashcodedark}
    sudo apt update
    sudo apt install -y ros-humble-turtlesim
\end{bashcodedark}

\vspace{0.3cm}

\Result{
    Después de instalarlo, podremos ejecutar el nodo principal del simulador.
}

\end{frame}

% =========================================================
\begin{frame}[fragile]{Ejecutar el simulador}

En una terminal, ejecute:

\begin{bashcodedark}
    ros2 run turtlesim turtlesim_node
\end{bashcodedark}

Deberá aparecer una ventana con una tortuga.

\vspace{0.3cm}

\DefBlock[ros2 run]{
    El comando \texttt{ros2 run} ejecuta un programa perteneciente a un paquete
    de ROS 2.
}

\end{frame}

% =========================================================
\begin{frame}[fragile]{Mover la tortuga con teclado}

Abra una segunda terminal y ejecute:

\begin{bashcodedark}
    ros2 run turtlesim turtle_teleop_key
\end{bashcodedark}

Use las flechas del teclado para mover la tortuga.

\vspace{0.3cm}

\Warning{
    La terminal donde se ejecuta \texttt{turtle\_teleop\_key} debe estar activa
    para que el teclado controle la tortuga.
}

\end{frame}

% =========================================================
\begin{frame}{¿Qué está ocurriendo?}

Cuando presionamos una tecla:

\begin{enumerate}
    \item el nodo de teleoperación genera un mensaje de velocidad;
    \item el mensaje se publica en \texttt{/turtle1/cmd\_vel};
    \item el nodo de Turtlesim recibe el mensaje;
    \item la tortuga modifica su movimiento;
    \item Turtlesim publica la nueva pose.
\end{enumerate}

\vspace{0.3cm}

\Result{
    Ya tenemos un sistema con nodos, tópicos y mensajes funcionando.
}

\end{frame}

% =========================================================
\section{Inspección de nodos}

\SectionSlide
{Inspección de nodos}
{Observar qué procesos están activos}

% =========================================================
\begin{frame}[fragile]{Listar nodos}

Con Turtlesim y teleop en ejecución, abra otra terminal y ejecute:

\begin{bashcodedark}
    ros2 node list
\end{bashcodedark}

Una salida posible es:

\begin{bashcode}
    /teleop_turtle
    /turtlesim
\end{bashcode}

\vspace{0.3cm}

\DefBlock[ros2 node list]{
    Muestra los nodos activos que actualmente participan en el sistema ROS 2.
}

\end{frame}

% =========================================================
\begin{frame}{Interpretación de los nodos}

\[
\texttt{/turtlesim}
\]

representa el nodo del simulador.

\[
\texttt{/teleop\_turtle}
\]

representa el nodo que lee el teclado y publica comandos de movimiento.

\vspace{0.3cm}

\begin{center}
    \begin{tikzpicture}[
        node distance=2.0cm,
        box/.style={
            ellipse,
            draw=BrandBlue,
            thick,
            rounded corners,
            minimum width=3.2cm,
            minimum height=0.9cm,
            align=center
        },
        arrow/.style={
            -{Latex[length=3mm]},
            thick,
            BrandBlue
        }
        ]

        \node[box] (teleop) {\texttt{/teleop\_turtle}};
        \node[box, right=of teleop] (sim) {\texttt{/turtlesim}};

        \draw[arrow] (teleop) -- node[above] {velocidad} (sim);

    \end{tikzpicture}
\end{center}

\end{frame}

% =========================================================
\section{Inspección de tópicos}

\SectionSlide
{Inspección de tópicos}
{Observar por dónde viajan los datos}

% =========================================================
\begin{frame}[fragile]{Listar tópicos}

Ejecute:

\begin{bashcodedark}
    ros2 topic list
\end{bashcodedark}

Entre los tópicos disponibles deberían aparecer:

\begin{bashcode}
    /turtle1/cmd_vel
    /turtle1/pose
\end{bashcode}

\vspace{0.3cm}

\DefBlock[ros2 topic list]{
    Muestra los canales de comunicación disponibles en el sistema ROS 2.
}

\end{frame}

% =========================================================
\begin{frame}{Dos tópicos importantes}

\[
\texttt{/turtle1/cmd\_vel}
\]

recibe comandos de velocidad.

\vspace{0.2cm}

\[
\texttt{/turtle1/pose}
\]

publica la posición y orientación de la tortuga.

\vspace{0.4cm}

\begin{center}
    \begin{tikzpicture}[
        node distance=1.2cm,
        box/.style={
            ellipse,
            draw=BrandBlue,
            thick,
            rounded corners,
            minimum width=3.2cm,
            minimum height=0.9cm,
            align=center
        },
        topic/.style={
            draw=BrandYellow,
            thick,
            rounded corners,
            minimum width=3.1cm,
            minimum height=0.85cm,
            align=center
        },
        arrow/.style={
            -{Latex[length=2.8mm]},
            thick,
            BrandBlue
        }
        ]

        \node[box] (cmdnode) {Nodo\\teleop};
        \node[topic, right=of cmdnode] (cmdvel) {\texttt{/cmd\_vel}};
        \node[box, right=of cmdvel] (sim) {Turtlesim};

        \node[topic, below=of sim] (pose) {\texttt{/pose}};

        \draw[arrow] (cmdnode) -- (cmdvel);
        \draw[arrow] (cmdvel) -- (sim);
        \draw[arrow] (sim) -- (pose);

    \end{tikzpicture}
\end{center}

\end{frame}

% =========================================================
\begin{frame}[fragile]{Información de un tópico}

Para conocer el tipo de mensaje de un tópico, ejecute:

\begin{bashcodedark}
    ros2 topic info /turtle1/cmd_vel
\end{bashcodedark}

Una salida posible es:

\begin{bashcode}
    Type: geometry_msgs/msg/Twist
    Publisher count: 1
    Subscription count: 1
\end{bashcode}

\vspace{0.3cm}

\Result{
    El tópico \texttt{/turtle1/cmd\_vel} transporta mensajes de tipo
    \texttt{Twist}.
}

\end{frame}

% =========================================================
\section{Mensaje Twist}

\SectionSlide
{Mensaje Twist}
{Velocidad lineal y velocidad angular}

% =========================================================
\begin{frame}{Estructura conceptual de Twist}

Un mensaje \texttt{Twist} contiene dos partes principales:

\[
\texttt{linear}
\]

y:

\[
\texttt{angular}
\]

\vspace{0.3cm}

\[
\texttt{linear.x}
\rightarrow
\text{velocidad lineal hacia adelante}
\]

\[
\texttt{angular.z}
\rightarrow
\text{velocidad angular de giro}
\]

\vspace{0.3cm}

\Idea{
    En Turtlesim trabajaremos principalmente con \texttt{linear.x}
    y \texttt{angular.z}.
}

\end{frame}

% =========================================================
\begin{frame}{Relación con el modelo matemático}

En robótica móvil plana solemos usar:

\[
v
\]

para la velocidad lineal y:

\[
\omega
\]

para la velocidad angular.

En ROS 2, para Turtlesim:

\[
v
\longrightarrow
\texttt{linear.x}
\]

\[
\omega
\longrightarrow
\texttt{angular.z}
\]

\vspace{-0.3cm}

\Result{
    La variable matemática se convierte en un campo dentro de un mensaje.
}

\end{frame}

% =========================================================
\begin{frame}{Interpretación del movimiento}

\begin{center}
    \begin{tikzpicture}[scale=0.9]

        % Ejes globales
        \draw[-{Latex[length=3mm]}, thick]
        (-0.5,0) -- (6,0)
        node[right] {$X$};

        \draw[-{Latex[length=3mm]}, thick]
        (0,-0.5) -- (0,4.5)
        node[above] {$Y$};

        % Posición del robot
        \coordinate (T) at (2.5,2);

        \begin{scope}[rotate around ={-90:(T)}]


        % Cuerpo circular
        \draw[
        fill=BrandBlue!15,
        draw=BrandBlue,
        thick
        ]
        (T) circle (0.55);

        % Rueda izquierda
        \draw[
        fill=black!70,
        draw=black,
        thick,
        rounded corners=1pt
        ]
        ($(T)+(-0.68,-0.23)$)
        rectangle
        ($(T)+(-0.53,0.23)$);

        % Rueda derecha
        \draw[
        fill=black!70,
        draw=black,
        thick,
        rounded corners=1pt
        ]
        ($(T)+(0.53,-0.23)$)
        rectangle
        ($(T)+(0.68,0.23)$);

        % Indicador del frente
        \draw[
        -{Latex[length=2.5mm]},
        very thick,
        BrandBlue
        ]
        (T) -- ++(0,0.42);

        % Velocidad lineal
        \draw[
        -{Latex[length=3mm]},
        very thick,
        BrandGreen
        ]
        ($(T)+(0,0.58)$) -- ++(0,1.2)
        node[above] {$v$};

        % Velocidad angular
        \draw[
        -{Latex[length=2.5mm]},
        BrandYellow,
        very thick
        ]
        ($(T)+(0.85,0)$)
        arc[
        start angle=0,
        end angle=120,
        radius=0.85
        ];

        \node[
        BrandYellow
        ]
        at ($(T)+(1.0,0.78)$)
        {$\omega$};

        % Centro del robot
        \fill[BrandBlue] (T) circle (1.7pt);

        % Etiqueta opcional
        \node[below=0.8cm] at (T) {robot};

        \end{scope}

    \end{tikzpicture}
\end{center}

\vspace{0.3cm}

\Idea{
    La velocidad lineal desplaza la tortuga. La velocidad angular cambia
    su orientación.
}

\end{frame}

% =========================================================
\section{Publicar mensajes desde terminal}

\SectionSlide
{Publicar mensajes}
{Mover Turtlesim sin teclado}

% =========================================================
\begin{frame}[fragile]{Publicar una velocidad}

Podemos enviar un mensaje directamente al tópico:

\begin{bashcodedark}
    ros2 topic pub /turtle1/cmd_vel geometry_msgs/msg/Twist \
    "{linear: {x: 2.0}, angular: {z: 0.0}}"
\end{bashcodedark}

Esto publica una velocidad lineal hacia adelante.

\vspace{0.3cm}

\Result{
    Si \texttt{angular.z} es cero, la tortuga avanza sin girar.
}

\end{frame}

% =========================================================
\begin{frame}[fragile]{Publicar con frecuencia limitada}

Para publicar un mensaje una sola vez:

\begin{bashcodedark}
    ros2 topic pub --once /turtle1/cmd_vel geometry_msgs/msg/Twist \
    "{linear: {x: 2.0}, angular: {z: 0.0}}"
\end{bashcodedark}

Para publicar a una frecuencia determinada:

\begin{bashcodedark}
    ros2 topic pub -r 2 /turtle1/cmd_vel geometry_msgs/msg/Twist \
    "{linear: {x: 1.0}, angular: {z: 0.5}}"
\end{bashcodedark}

\vspace{0.3cm}

\Warning{
    Un comando con \texttt{-r} sigue publicando hasta detenerlo con
    \texttt{Ctrl + C}.
}

\end{frame}

% =========================================================
\begin{frame}{Efecto de los comandos}

\begin{center}
    \begin{tabular}{ccc}
        \toprule
        \textbf{\texttt{linear.x}} & \textbf{\texttt{angular.z}} & \textbf{Efecto} \\
        \midrule
        $>0$ & $0$ & avanza recto \\
        $0$ & $>0$ & gira antihorario \\
        $0$ & $<0$ & gira horario \\
        $>0$ & $>0$ & curva antihoraria \\
        $>0$ & $<0$ & curva horaria \\
        \bottomrule
    \end{tabular}
\end{center}

\vspace{0.3cm}

\Idea{
    El movimiento de la tortuga permite observar directamente la relación
    entre velocidad lineal y velocidad angular.
}

\end{frame}

% =========================================================
\section{Observar la pose}

\SectionSlide
{Observar la pose}
{Posición y orientación publicadas por Turtlesim}

% =========================================================
\begin{frame}[fragile]{Escuchar el tópico de pose}

Ejecute:

\begin{bashcodedark}
    ros2 topic echo /turtle1/pose
\end{bashcodedark}

Mientras la tortuga se mueve, observe cómo cambian:

\begin{itemize}
    \item \texttt{x};
    \item \texttt{y};
    \item \texttt{theta};
    \item \texttt{linear\_velocity};
    \item \texttt{angular\_velocity}.
\end{itemize}

\end{frame}

% =========================================================
\begin{frame}{Interpretación de la pose}

    La pose de Turtlesim contiene:

    \[
    \begin{bmatrix}
        x \rightarrow \text{posición horizontal.} \\

        y \rightarrow \text{posición vertical.} \\


        \theta \rightarrow \text{orientación de la tortuga.}\\

    \end{bmatrix}
    \]

    \[
    \mathbf{x}
    =
    \begin{bmatrix}
        x\\
        y\\
        \theta
    \end{bmatrix}
    \]

    \vspace{-0.3cm}

    \DefBlock[Pose]{
        La pose reúne posición y orientación. En robótica móvil plana suele
        representarse como $(x,y,\theta)$.
    }

\end{frame}

% =========================================================
\begin{frame}{Pose y orientación}

\begin{center}
    \begin{tikzpicture}[scale=0.85]

        % Ejes globales
        \draw[-{Latex[length=3mm]}, thick]
        (0,0) -- (6.5,0)
        node[right] {$X$};

        \draw[-{Latex[length=3mm]}, thick]
        (0,0) -- (0,5)
        node[above] {$Y$};

        % Posición de la tortuga
        \coordinate (R) at (3.5,2.4);

        % Proyecciones de posición
        \draw[dashed, gray]
        (R) -- (3.5,0);

        \draw[dashed, gray]
        (R) -- (0,2.4);

        % Coordenadas
        \node[below] at (3.5,0) {$x$};
        \node[left]  at (0,2.4) {$y$};

        % Tortuga rotada según theta
        \begin{scope}[rotate around={35:(R)}]

            % Caparazón
            \draw[
            fill=BrandBlue!15,
            draw=BrandBlue,
            very thick
            ]
            (R) ellipse (0.62 and 0.42);

            % Cabeza
            \draw[
            fill=BrandBlue!25,
            draw=BrandBlue,
            thick
            ]
            ($(R)+(0.78,0)$) circle (0.18);

            % Patas
            \draw[
            fill=BrandBlue!25,
            draw=BrandBlue,
            thick
            ]
            ($(R)+(0.25,0.48)$) ellipse (0.17 and 0.10);

            \draw[
            fill=BrandBlue!25,
            draw=BrandBlue,
            thick
            ]
            ($(R)+(0.25,-0.48)$) ellipse (0.17 and 0.10);

            \draw[
            fill=BrandBlue!25,
            draw=BrandBlue,
            thick
            ]
            ($(R)+(-0.30,0.48)$) ellipse (0.17 and 0.10);

            \draw[
            fill=BrandBlue!25,
            draw=BrandBlue,
            thick
            ]
            ($(R)+(-0.30,-0.48)$) ellipse (0.17 and 0.10);

            % Cola
            \draw[
            fill=BrandBlue!25,
            draw=BrandBlue,
            thick
            ]
            ($(R)+(-0.70,0)$) -- ++(-0.22,0.10)
            -- ++(0,-0.20) -- cycle;

        \end{scope}

        % Centro de referencia
        \fill[BrandBlue] (R) circle (2pt);

        % Dirección de la tortuga
        \draw[
        -{Latex[length=3mm]},
        very thick,
        BrandGreen
        ]
        (R) -- ++(35:1.8)
        node[right] {dirección};

        % Ángulo theta
        \draw[
        -{Latex[length=2.5mm]},
        BrandYellow,
        very thick
        ]
        ($(R)+(1.0,0)$)
        arc[
        start angle=0,
        end angle=35,
        radius=0.9
        ];

        \node[BrandYellow]
        at ($(R)+(1.2,0.3)$)
        {$\theta$};

        % Etiqueta
        \node[above left=0.45cm] at (R) {Tortuga};

    \end{tikzpicture}
\end{center}

\vspace{0.2cm}

\Result{
    El tópico \texttt{/turtle1/pose} permite observar el estado geométrico
    de la tortuga.
}

\end{frame}

% =========================================================
\section{Tipos de mensaje}

\SectionSlide
{Tipos de mensaje}
{Saber qué estructura viaja por un tópico}

% =========================================================
\begin{frame}[fragile]{Ver el tipo de mensaje}

Para conocer el tipo de mensaje de la pose:

\begin{bashcodedark}
    ros2 topic info /turtle1/pose
\end{bashcodedark}

Una salida posible es:

\begin{bashcode}
    Type: turtlesim/msg/Pose
    Publisher count: 1
    Subscription count: 1
\end{bashcode}

\vspace{0.3cm}

\Result{
    El tópico \texttt{/turtle1/pose} utiliza mensajes de tipo
    \texttt{turtlesim/msg/Pose}.
}

\end{frame}

% =========================================================
\begin{frame}[fragile]{Ver la estructura de un mensaje}

Ejecute:

\begin{bashcodedark}
    ros2 interface show turtlesim/msg/Pose
\end{bashcodedark}

Una salida posible es:

\begin{bashcode}
    float32 x
    float32 y
    float32 theta
    float32 linear_velocity
    float32 angular_velocity
\end{bashcode}

\vspace{0.3cm}

\Idea{
    La interfaz del mensaje indica qué datos contiene y cómo están organizados.
}

\end{frame}

% =========================================================
\begin{frame}[fragile]{Estructura de Twist}

Para ver la estructura de \texttt{Twist}:

\begin{bashcodedark}
    ros2 interface show geometry_msgs/msg/Twist
\end{bashcodedark}

La estructura general es:

\begin{bashcode}
    Vector3  linear
    Vector3  angular
\end{bashcode}

\vspace{0.3cm}

\DefBlock[Twist]{
    Un mensaje \texttt{Twist} representa velocidad lineal y velocidad angular.
}

\end{frame}

% =========================================================
\section{Laboratorio guiado}

\SectionSlide
{Laboratorio guiado}
{Explorar nodos, tópicos, mensajes y movimiento}

% =========================================================
\begin{frame}{Parte 1: iniciar el sistema}

Realice los siguientes pasos:

\begin{enumerate}
    \item abra una terminal;
    \item cargue ROS 2 Humble si es necesario;
    \item ejecute \texttt{turtlesim\_node};
    \item abra una segunda terminal;
    \item ejecute \texttt{turtle\_teleop\_key};
    \item mueva la tortuga con el teclado.
\end{enumerate}

\vspace{0.3cm}

\Practice{
    Observe qué ocurre cuando se presionan las flechas y relacione el
    movimiento con velocidad lineal y angular.
}

\end{frame}

% =========================================================
\begin{frame}{Parte 2: inspección}

Con el sistema en ejecución:

\begin{enumerate}
    \item liste los nodos activos;
    \item liste los tópicos disponibles;
    \item consulte la información de \texttt{/turtle1/cmd\_vel};
    \item consulte la información de \texttt{/turtle1/pose};
    \item observe la estructura de \texttt{geometry\_msgs/msg/Twist};
    \item observe la estructura de \texttt{turtlesim/msg/Pose}.
\end{enumerate}

\vspace{0.3cm}

\Result{
    Esta parte permite relacionar lo que se ve en pantalla con la arquitectura
    interna de ROS 2.
}

\end{frame}

% =========================================================
\begin{frame}{Parte 3: publicación manual}

Publique diferentes comandos de velocidad.

Pruebe:

\begin{itemize}
    \item avance recto;
    \item giro en el sitio;
    \item curva hacia la izquierda;
    \item curva hacia la derecha;
    \item detener la tortuga.
\end{itemize}

\vspace{0.3cm}

\Idea{
    Cada movimiento debe interpretarse según los valores asignados a
    \texttt{linear.x} y \texttt{angular.z}.
}

\end{frame}

% =========================================================
\begin{frame}[fragile]{Comandos sugeridos}

Avanzar recto:

\begin{bashcodedark}
    ros2 topic pub --once /turtle1/cmd_vel geometry_msgs/msg/Twist \
    "{linear: {x: 2.0}, angular: {z: 0.0}}"
\end{bashcodedark}

Girar en el sitio:

\begin{bashcodedark}
    ros2 topic pub --once /turtle1/cmd_vel geometry_msgs/msg/Twist \
    "{linear: {x: 0.0}, angular: {z: 1.5}}"
\end{bashcodedark}

\end{frame}

% =========================================================
\begin{frame}[fragile]{Comandos sugeridos}

Curva antihoraria:

\begin{bashcodedark}
    ros2 topic pub -r 2 /turtle1/cmd_vel geometry_msgs/msg/Twist \
    "{linear: {x: 1.0}, angular: {z: 0.7}}"
\end{bashcodedark}

Curva horaria:

\begin{bashcodedark}
    ros2 topic pub -r 2 /turtle1/cmd_vel geometry_msgs/msg/Twist \
    "{linear: {x: 1.0}, angular: {z: -0.7}}"
\end{bashcodedark}

Detenga con:

\begin{center}
    \texttt{Ctrl + C}
\end{center}

\end{frame}

% =========================================================
\begin{frame}{Parte 4: observar la pose}

Mientras la tortuga se mueve:

\begin{enumerate}
    \item observe \texttt{/turtle1/pose};
    \item registre tres valores diferentes de $(x,y,\theta)$;
    \item identifique cuándo cambia más \texttt{x};
    \item identifique cuándo cambia más \texttt{y};
    \item observe qué ocurre con \texttt{theta} durante un giro.
\end{enumerate}

\vspace{0.3cm}

\Practice{
    Relacione el valor de \texttt{theta} con la orientación geométrica
    de la tortuga en la ventana.
}

\end{frame}

% =========================================================
\section{Cierre}

\SectionSlide
{Cierre de la sesión}
{ROS 2 como infraestructura de comunicación}

% =========================================================
\begin{frame}{¿Qué construimos hoy?}

Hoy trabajamos con:

\begin{itemize}
    \item ROS 2 Humble;
    \item Turtlesim;
    \item nodos;
    \item tópicos;
    \item mensajes;
    \item publicadores;
    \item suscriptores;
    \item comandos básicos de inspección;
    \item mensajes \texttt{Twist};
    \item pose de la tortuga.
\end{itemize}

\end{frame}

% =========================================================
\begin{frame}{La idea más importante}

\[
\boxed{
    \text{Un robot en ROS 2 es un conjunto de nodos que intercambian mensajes}
}
\]

\vspace{0.4cm}

\[
\text{Nodo}
\rightarrow
\text{Tópico}
\rightarrow
\text{Mensaje}
\rightarrow
\text{Acción}
\]

\vspace{0.3cm}

\Idea{
    Antes de programar nodos propios, es necesario comprender cómo observar,
    publicar e interpretar información usando la terminal.
}

\end{frame}

% =========================================================
\begin{frame}{Preguntas de cierre}

Responda:

\begin{enumerate}
    \item ¿Qué diferencia existe entre un nodo y un tópico?
    \item ¿Qué función cumple un mensaje?
    \item ¿Por qué \texttt{/turtle1/cmd\_vel} usa un mensaje \texttt{Twist}?
    \item ¿Qué representa \texttt{linear.x}?
    \item ¿Qué representa \texttt{angular.z}?
    \item ¿Qué información contiene \texttt{/turtle1/pose}?
\end{enumerate}

\end{frame}

% =========================================================
\begin{frame}{Próxima clase}

En la siguiente sesión trabajaremos con programación en Python para ROS 2.

El objetivo será construir nuestros primeros nodos usando:

\[
\boxed{\texttt{rclpy}}
\]

y desarrollar:

\begin{itemize}
    \item un publicador de velocidad;
    \item un suscriptor de pose;
    \item un movimiento controlado de Turtlesim;
    \item una primera relación entre código Python y tópicos ROS 2.
\end{itemize}

\end{frame}

% =========================================================
\begin{frame}
\centering
\smartimage[0.8\linewidth]{Template/gracias.png}
\end{frame}

\end{document}